\chapter{Module 5 : notifications, analytique, administration et gouvernance}

\section*{Introduction}
\addcontentsline{toc}{section}{Introduction}
Le dernier module referme la boucle fonctionnelle de la plateforme. Après avoir construit l'accès, le vivier documentaire, le pipeline candidat et la couche IA, il restait à renforcer la gouvernance globale du système. Une plateforme de recrutement d'entreprise ne peut pas se limiter à traiter des candidatures ; elle doit aussi fournir de la visibilité, des traces d'audit, des mécanismes de communication, des exports et un suivi après embauche.

Ce module ne correspond donc pas à un simple bloc de notifications. Il consolide la dimension administrable, pilotable et auditable de \projectname{}, là où la supervision, la traçabilité et la restitution des données deviennent aussi importantes que les fonctionnalités de recrutement elles-mêmes.

Il montre ainsi comment le produit passe d'un système intelligent à un système gouvernable. Les données, décisions et automatisations produites précédemment doivent désormais pouvoir être suivies, expliquées, contrôlées et exportées.

\section{Sprint 5 : objectif et périmètre}

\subsection{Objectif du sprint}
L'objectif est de donner aux utilisateurs et aux administrateurs les moyens de superviser le système, de suivre les événements importants et de prolonger le recrutement jusqu'à l'intégration effective du candidat recruté. Plus précisément, ce sprint vise à :
\begin{itemize}
  \item notifier les acteurs au bon moment ;
  \item conserver la mémoire des actions, des emails et des changements de statut ;
  \item fournir aux administrateurs une lecture consolidée de la plateforme ;
  \item offrir des exports exploitables pour le suivi et l'audit ;
  \item prolonger le cycle de vie du recrutement jusqu'à l'onboarding.
\end{itemize}

Il répond à une contrainte réaliste des outils d'entreprise : une bonne automatisation n'a de valeur durable que si elle laisse des preuves, des historiques et des points de pilotage lisibles pour les équipes responsables.

\subsection{Backlog du sprint}
\begin{table}[H]
  \centering
  \caption{Backlog fonctionnel du module gouvernance et administration}
  \begin{tabularx}{\textwidth}{>{\bfseries}p{1.3cm}p{3cm}X>{\raggedright\arraybackslash}p{2cm}}
    \toprule
    ID & Thème & User Story & Statut \\
    \midrule
    33 & Notifications & En tant qu'utilisateur connecté, je peux recevoir et consulter des notifications in-app. & Réalisé \\
    34 & Rappels d'entretien & En tant qu'utilisateur concerné, je peux recevoir des rappels d'entretien du jour. & Réalisé \\
    35 & Logs d'emails & En tant qu'administrateur, je peux consulter l'historique des emails envoyés. & Réalisé \\
    36 & Journal d'activité & En tant qu'administrateur, je peux auditer les actions effectuées sur la plateforme. & Réalisé \\
    37 & Vue système & En tant qu'administrateur, je peux consulter un aperçu global de la plateforme. & Réalisé \\
    38 & Analytics RH & En tant qu'administrateur, je peux analyser les performances du recrutement et la santé du pipeline. & Réalisé \\
    39 & Exports & En tant qu'administrateur, je peux exporter les données clés en Excel ou ZIP. & Réalisé \\
    40 & Onboarding détaillé & En tant qu'administrateur, je peux suivre l'avancement d'intégration des candidats embauchés. & Réalisé \\
    41 & Centre de gouvernance & En tant qu'administrateur, je peux filtrer et exporter les transitions candidat, les confirmations IA et les logs d'activité depuis une même vue. & Réalisé \\
    42 & Durcissement release & En tant qu'équipe projet, je peux vérifier que les migrations critiques et les comportements d'audit sont présents avant livraison. & Réalisé \\
    43 & Qualité agentique & En tant qu'équipe projet, nous pouvons valider les cartes de réponse, les traces d'outils, les confirmations et le masquage des données sensibles. & Réalisé \\
    \bottomrule
  \end{tabularx}
\end{table}

Ce backlog confirme que la gouvernance n'est pas un ajout marginal. Elle couvre la communication, l'audit, l'analytique, l'export, les confirmations IA, le centre de gouvernance et la continuité post-recrutement.

\section{Implémentation du sprint}

\subsection{Analyse du sprint}
Ce sprint s'appuie principalement sur les éléments fonctionnels suivants :
\begin{itemize}
  \item le service de notifications in-app ;
  \item le service de communications sortantes et de journalisation des emails ;
  \item le service de vues agrégées administrateur ;
  \item le service de traçabilité métier ;
  \item le service d'export des données ;
  \item le service de checklists d'intégration ;
  \item le service de rapport d'audit gouvernance ;
  \item l'interface de notification ;
  \item la vue système administrateur ;
  \item la vue analytique administrateur ;
  \item la vue du journal d'activité ;
  \item la vue des logs d'emails ;
  \item la vue de suivi détaillé de l'intégration ;
  \item le centre de gouvernance administrateur ;
  \item les tests de qualité agentique couvrant confirmations, preuves, cartes et traces.
\end{itemize}

Ce sprint crée donc une couche de gouvernance transversale qui s'appuie sur tout ce qui a été produit auparavant : candidats, entretiens, décisions, emails, IA et onboarding.

Cette transversalité montre que la supervision n'est pas plaquée après coup. Elle réutilise les mêmes données opérationnelles pour produire une lecture de pilotage et une capacité d'audit.

\subsection{Vue fonctionnelle globale}
Les usages couverts se regroupent en cinq blocs :
\begin{enumerate}
  \item \textbf{Communication et alertes} : créer des notifications, les récupérer, compter les éléments non lus, marquer une notification comme lue, tout marquer comme lu et envoyer ou tracer les emails.
  \item \textbf{Supervision administrateur} : consulter un aperçu global du système, analyser les métriques du recrutement et suivre l'activité transversale de la plateforme.
  \item \textbf{Restitution et audit} : conserver un journal enrichi des actions, un historique des emails et exporter les données clés au format Excel ou ZIP.
  \item \textbf{Continuité post-recrutement} : suivre les checklists d'onboarding, visualiser l'état d'intégration des candidats \sourcecode{hired} et prolonger le pilotage au-delà de la décision d'embauche.
  \item \textbf{Gouvernance des actions sensibles} : rapprocher les transitions candidat, les confirmations d'actions IA et les logs d'activité dans une vue filtrable et exportable.
\end{enumerate}

Cette structuration illustre la maturité atteinte par le produit : autour des entités métier centrales s'ajoutent désormais des objets de preuve, de communication et de suivi.

\subsection{Chaîne de notification et de communication}
Le service \sourcecode{notifications.ts} gère les notifications en base via la table \sourcecode{notifications}. Le système peut créer une notification ciblée, récupérer les dernières notifications d'un utilisateur, calculer le nombre non lu, marquer une notification comme lue et tout marquer comme lu.

Deux cas métier sont particulièrement importants : \sourcecode{notifyStageChange()} pour les changements d'étape d'un candidat et \sourcecode{notifyInterviewScheduled()} pour les événements liés à la planification d'entretien. En complément, \sourcecode{ensureTodayInterviewReminders()} vérifie les entretiens du jour afin de générer les rappels nécessaires.

Le service \sourcecode{email.ts} ne se contente pas d'envoyer des messages. Il trace aussi ces envois dans \sourcecode{email\_logs}, ce qui garantit qu'une communication importante du recrutement reste auditée et restituable.

Cette couche améliore la réactivité du produit et réduit le risque qu'une étape importante reste invisible, en particulier dans un processus où plusieurs rôles interviennent sur des moments différents.

\begin{figure}[H]
  \centering
  \diagraminput{figures/diagrams/notification-email-sequence}
  \caption{Chaîne de notification et de journalisation des emails}
\end{figure}

\subsection{Supervision administrateur}
La couche \sourcecode{admin.ts} agrège plusieurs vues de haut niveau : \sourcecode{getSystemOverview()} calcule une vision consolidée du système, \sourcecode{getRecruitmentAnalytics()} fournit des métriques de performance et de répartition, \sourcecode{getEmailLogs()} expose l'historique des emails et \sourcecode{getHiredCandidatesOnboardingDetailed()} restitue l'état détaillé de l'onboarding.
Cette agrégation montre que la plateforme ne sépare pas l'opérationnel du décisionnel. Les données utilisées pour exécuter le recrutement sont aussi réutilisées pour le piloter.
La version actuelle ajoute à ces surfaces des panneaux d'évidence administrateur et des entrées agentiques gouvernées. Chaque vue peut maintenant résumer les faits observés, signaler les données opérationnelles manquantes et proposer des actions \textit{Agent} dédiées à la synthèse d'activité, à l'interprétation analytique, aux anomalies d'onboarding ou aux risques de communication.
Cette réutilisation des mêmes données constitue un point fort du projet : elle évite de dupliquer la vérité métier dans un outil de reporting distinct et renforce la cohérence entre exécution, supervision et assistance à la décision.
Le nouveau centre de gouvernance renforce cette lecture. Il ne se limite pas au journal d'activité classique : il agrège les transitions du pipeline, les actions IA en attente ou exécutées et les logs applicatifs, afin de fournir une surface de contrôle unique pour les événements sensibles.
\paragraph{Lecture de gouvernance et portée de supervision.}

Elle permet surtout de distinguer clairement deux niveaux de lecture. Le niveau opérationnel sert à faire avancer les dossiers candidats, à planifier des entretiens et à envoyer des communications. Le niveau de gouvernance sert, lui, à vérifier que ces actions ont bien été réalisées, à en mesurer les effets, à en conserver la preuve et à fournir aux responsables une vision consolidée du système.


\subsection{Audit, logs et exports}
Le service \sourcecode{activity-log.ts} enregistre les événements importants via \sourcecode{logActivity()}, puis les restitue en lecture simple ou enrichie. En parallèle, \sourcecode{export.ts} permet de générer plusieurs artefacts : export du vivier de CV, export d'un CV individuel, export ZIP de plusieurs CV, export des candidats acceptés, export des logs d'emails, export du journal d'activité, export de l'onboarding et export Excel détaillé d'un workflow d'acceptation candidat.

Cette partie prouve que la plateforme est pensée pour la restitution, la preuve et l'audit, pas seulement pour l'exécution interactive. Les exports jouent un rôle d'interopérabilité : la donnée reste exploitable hors de l'interface, pour le reporting, le contrôle et la conformité.
Dans un cadre d'entreprise, cette capacité d'export dépasse la simple commodité. Elle permet de transformer l'information applicative en matière de reporting, de contrôle, de conformité et de partage. Les exports constituent donc un pont entre la plateforme de recrutement et les besoins plus larges d'audit, de pilotage RH ou de restitution managériale.

\subsection{Centre de gouvernance et audit consolidé}
Le service \sourcecode{governance.ts} construit un rapport d'audit consolidé à partir de trois sources : \sourcecode{candidate\_stage\_history}, \sourcecode{pending\_agent\_actions} et \sourcecode{activity\_logs}. Chaque ligne est normalisée avec un type, un statut, une action, un acteur, un candidat éventuel, un horodatage et un détail métier.

Cette consolidation répond à un besoin de contrôle que les vues séparées ne couvrent pas entièrement. Un administrateur peut filtrer les événements par période, acteur, candidat, action ou statut, puis exporter les lignes en CSV. Les arguments d'outils sont sanitizés avant restitution afin d'éviter d'exposer des données sensibles comme les tokens, les pièces jointes brutes ou les contenus volumineux.

Le centre de gouvernance joue donc le rôle d'un plan de contrôle source-backed. Il met au même niveau les changements d'étape, les actions agentiques sensibles et les logs applicatifs, ce qui facilite l'analyse d'incident, la revue d'une décision et la préparation d'un audit.


\section{Réalisation}

\subsection{Cloche de notification et communication in-app}
Le composant \sourcecode{notification-bell.tsx}, injecté dans \sourcecode{DashboardShell}, matérialise la couche de notification côté interface. Il permet de charger les notifications à l'ouverture du panneau, d'afficher le compteur d'éléments non lus, de marquer une notification individuellement et de tout marquer comme lu. Cette intégration rend la gouvernance visible dans l'expérience quotidienne.

Elle renforce aussi la dimension temps réel du produit : les événements importants remontent au plus près du point d'usage, au lieu de rester enfouis dans des vues administratives.

\begin{figure}[H]
  \centering
  \screenshotplaceholder{Notification bell/state}
  \caption{Cloche de notification et état de lecture.}
\end{figure}

\subsection{Tableau de bord administrateur}
La page de tableau de bord administrateur appelle \sourcecode{getSystemOverviewAction()} et alimente \sourcecode{AdminDashboardClient}. Cette page donne une vision d'ensemble sur l'état du produit avec un discours orienté supervision : volumes, activité et santé fonctionnelle du pipeline.
Elle combine maintenant les cartes KPI avec un panneau \textit{Governance snapshot} source-backed et plusieurs CTAs agentiques dédiés à la synthèse d'activité, à la revue des rôles, aux anomalies opérationnelles et au brief de gouvernance.
C'est le point de contrôle macro du système, là où l'administrateur peut comprendre rapidement la situation globale de la plateforme puis basculer, si nécessaire, vers une analyse agentique plus approfondie.

\begin{figure}[H]
  \centering
  \screenshotplaceholder{Admin dashboard}
  \caption{Tableau de bord administrateur.}
\end{figure}

\subsection{Analytique administrateur}
La page analytique administrateur repose sur \sourcecode{getRecruitmentAnalyticsAction()} et sur \sourcecode{AdminAnalyticsClient}. Elle transforme les données opérationnelles en indicateurs lisibles : performance du pipeline, répartition des volumes, tendances de recrutement et signaux utiles au pilotage global.
Elle ne se limite plus aux graphiques. Un panneau d'évidence dédié récapitule désormais les faits observés, les limites de mesure et les risques d'interprétation, tandis que des actions \textit{Explain trends}, \textit{Find bottlenecks} et \textit{Governance risks} ouvrent l'agent avec un contexte déjà cadré.
La restitution analytique combine ainsi deux niveaux complémentaires : des graphiques de page pour la supervision globale et des réponses agentiques capables de produire des visualisations ciblées lorsque l'utilisateur demande une courbe, une répartition ou une comparaison demande--offre.
Elle donne ainsi à l'administration la possibilité de lire les flux de recrutement comme un ensemble dynamique, et non comme une simple accumulation de dossiers.

\begin{figure}[H]
  \centering
  \screenshotplaceholder{Admin analytics}
  \caption{Vue analytique administrateur.}
\end{figure}

\subsection{Centre de gouvernance}
La page \sourcecode{/admin/governance} charge \sourcecode{getGovernanceAuditReportAction()} puis affiche \sourcecode{AdminGovernanceClient}. Elle présente des compteurs dédiés aux lignes d'audit, aux changements d'étape, aux actions IA, aux logs d'activité, aux actions IA en attente et aux actions IA échouées.

L'interface permet de filtrer par date, acteur, candidat, action ou statut. Chaque ligne ouvre un panneau de détail contenant une preuve sanitizée, avec les arguments d'outil nettoyés lorsque l'événement provient de l'agent. L'administrateur peut aussi exporter le résultat courant en CSV via \sourcecode{exportGovernanceAuditCsvAction()}.

\begin{figure}[H]
  \centering
  \screenshotplaceholder{Governance Center}
  \caption{Centre de gouvernance administrateur.}
\end{figure}

\subsection{Historique d'activité et logs d'emails}
La page du journal d'activité exploite \sourcecode{getActivityLogEnrichedAction(100)} et le composant \sourcecode{AdminActivityClient}. L'objectif est d'exposer une lecture enrichie du journal d'activité afin de reconstituer les séquences d'action importantes. Un panneau d'évidence et des actions \textit{Agent} permettent en parallèle de synthétiser les lignes chargées, de signaler les actions destructrices et de rappeler les limites de couverture de l'audit courant.
Cette fonctionnalité est centrale pour l'audit interne, l'analyse d'incident, la compréhension des décisions et la justification d'un enchaînement métier. Dans une plateforme de recrutement augmentée par l'IA, cette traçabilité a encore plus de valeur, car elle permet d'observer ce qui a été automatisé, ce qui a été décidé et dans quel ordre.
Le journal d'activité ne sert donc pas seulement à voir ce qui s'est passé. Il sert aussi à rendre le système explicable lorsqu'une décision, une communication ou un changement de statut doit être revu.
\begin{figure}[H]
  \centering
  \screenshotplaceholder{Activity log}
  \caption{Journal d'activité.}
\end{figure}
La page des logs d'emails charge \sourcecode{getEmailLogsAction()} puis transmet les données au composant \sourcecode{AdminEmailsClient}. Cette vue constitue l'\textit{audit trail} des communications sortantes de la plateforme, mais elle bénéficie à présent des mêmes principes : panneau \textit{Communication evidence readiness}, CTAs agentiques dédiés aux risques de livraison et liens profonds vers les entrées auditées.
Elle est particulièrement importante dans le contexte RH, où les invitations, confirmations, refus et communications de décision doivent rester consultables et vérifiables. La consultation centralisée de ces logs évite qu'une partie sensible du processus reste cachée dans une boîte noire extérieure au produit.
Cette séparation entre journal d'activité et logs d'emails clarifie la gouvernance : le premier explique les actions métier réalisées dans l'application, tandis que les seconds documentent les communications envoyées vers les candidats ou les parties prenantes.

\subsection{Suivi détaillé de l'onboarding}
La page de suivi détaillé de l'onboarding s'appuie sur \sourcecode{getHiredCandidatesOnboardingDetailedAction()} et \sourcecode{AdminOnboardingClient}. Elle ne montre pas seulement une liste de candidats embauchés ; elle expose aussi l'état détaillé des tâches d'intégration, ce qui prolonge le recrutement après la décision finale.
Le tableau embarque désormais un panneau d'évidence consacré aux anomalies d'onboarding : checklists absentes, coordonnées incomplètes, profils pauvres en signaux et autres trous de couverture. Les CTAs \textit{Find anomalies}, \textit{Summarize readiness} et \textit{Risk review} transmettent ensuite ces constats à l'agent sans réintroduire de génération libre.
Ce prolongement vers l'onboarding est important pour la cohérence globale du produit : le cycle de vie géré par la plateforme ne s'arrête pas à la validation finale, mais continue jusqu'au suivi de l'entrée effective dans l'organisation.

\begin{figure}[H]
  \centering
  \screenshotplaceholder{Onboarding tracker}
  \caption{Suivi de l'onboarding.}
\end{figure}

\subsection{Exports et interopérabilité}
Les actions d'export disponibles dans \sourcecode{actions.ts} et \sourcecode{export.ts} donnent à la plateforme une capacité d'interopérabilité concrète. Les données peuvent sortir sous des formats directement exploitables par les équipes métier, de contrôle ou d'audit. Le produit n'enferme donc pas l'information dans son interface.

Cette capacité d'export renforce la maturité du système : un produit d'entreprise crédible doit aussi pouvoir alimenter d'autres usages organisationnels hors de son interface principale.

\subsection{Vérification de durcissement release}
Le script \sourcecode{db/verify-release-hardening.ts} complète la gouvernance par une vérification technique ciblée. Il contrôle la présence des tables et index critiques liés à l'historique des étapes et aux confirmations IA, peut appliquer les migrations manquantes lorsque le schéma est absent, puis exécute un scénario comportemental temporaire.

Ce scénario crée un utilisateur de vérification, un poste, un CV et un candidat, valide que l'affectation et une transition manuelle produisent bien l'historique attendu, puis confirme une action IA qui fait progresser le candidat. Le script vérifie ensuite que l'action passe au statut \sourcecode{executed}, que la transition agentique est enregistrée, puis supprime les lignes temporaires. Il ne s'agit donc pas d'un simple contrôle de schéma, mais d'une preuve de bout en bout sur les mécanismes de gouvernance les plus sensibles.

Cette logique est complétée par les tests automatisés de l'agent et de son interface. Ils vérifient que les actions sensibles restent en attente de confirmation, que les cartes candidat/pipeline/gouvernance proviennent de résultats d'outils, que les traces d'exécution sont classées par phases et que les champs sensibles sont masqués avant affichage ou stockage technique. Cette couche de vérification relie donc directement l'expérience utilisateur du chat aux exigences de gouvernance décrites dans ce chapitre.

\paragraph{Lecture d'audit.}
La gouvernance du module repose sur une séparation claire des traces. Les notifications servent à informer les utilisateurs au moment utile. Le journal d'activité permet de reconstituer les actions réalisées dans l'application. Les logs d'emails documentent les communications sortantes. Les exports rendent ces informations exploitables hors de l'interface. Cette séparation évite de mélanger alerte, preuve, communication et restitution, ce qui renforce la lisibilité du système en cas de contrôle ou de revue métier.

\section{Apport du module}
Ce module apporte la maturité de gouvernance du projet :
\begin{itemize}
  \item il améliore la réactivité via les notifications et les rappels ;
  \item il améliore la traçabilité via les logs d'activité et d'emails ;
  \item il améliore le pilotage grâce aux dashboards et aux analytics administrateurs ;
  \item il améliore l'interopérabilité par les exports ;
  \item il prolonge la continuité métier jusqu'à l'onboarding ;
  \item il améliore la gouvernance des actions IA sensibles par une confirmation explicite et auditée ;
  \item il fournit un centre de gouvernance consolidé pour les transitions, confirmations et logs ;
  \item il renforce la qualité agentique par des tests couvrant confirmations, cartes typées, traces par phases et masquage des données sensibles.
\end{itemize}

Dans l'économie générale du rapport, ce module montre que l'intelligence et l'automatisation déployées dans les chapitres précédents sont encadrées par une couche de supervision explicite. C'est cette combinaison entre IA, métier et gouvernance qui donne à \projectname{} sa crédibilité comme produit d'entreprise.

On peut donc considérer ce dernier module comme la couche de confiance et de pilotage finale du système. Il transforme la plateforme en produit non seulement fonctionnel et intelligent, mais aussi administrable, explicable et durable dans un contexte professionnel.

\section*{Conclusion}
\addcontentsline{toc}{section}{Conclusion}
Ce cinquième module transforme la plateforme en système pilotable, gouvernable et auditable. Notifications, logs, dashboards administrateurs, analytics, centre de gouvernance, confirmations IA, exports et suivi d'onboarding complètent désormais les modules métier et IA déjà mis en place. Le produit ne se limite plus à exécuter un recrutement intelligent ; il permet aussi de le superviser, de l'expliquer et de le prolonger au-delà de la décision finale.
